\documentclass[preprint,12pt]{elsarticle}
\usepackage[T1]{fontenc}
\usepackage[utf8]{inputenc}
\usepackage[brazil]{babel}
\usepackage{amsmath,amssymb,booktabs,graphicx,xurl,placeins}
\usepackage[hidelinks]{hyperref}
\setlength{\emergencystretch}{2em}
\renewcommand{\floatpagefraction}{0.8}
\renewcommand{\topfraction}{0.9}
\renewcommand{\textfraction}{0.1}
\journal{Supply Chain Analytics -- manuscrito de trabalho em português}
\begin{document}
\begin{frontmatter}
\title{AIPE: seleção contextual de políticas de reposição por perfil e por aprendizado supervisionado}
\author[ufal]{Matheus Inácio Batista Santos}
\affiliation[ufal]{organization={Instituto de Computação, Universidade Federal de Alagoas},country={Brasil}}
\begin{abstract}
Escolher uma política de reposição exige relacionar a demanda às consequências de custo e atendimento. Apresentamos o Adaptive Inventory Policy Engine (AIPE), que integra caracterização, calibração de candidatas e seleção contextual. Seu Policy Selection Engine (PSE) admite associação entre perfil e política e previsão supervisionada da política por série. Avaliamos sete candidatas com separação temporal e seleção em outras lojas. Na Bahia, 145 séries em duas origens totalizam 4.350 ajustes; o PSE por perfil aumenta o atendimento em 3,89 pontos percentuais e o custo em 1,93\% frente à seleção global. Em um piloto M5 com 30 séries, dez lojas e três origens, os 1.350 ajustes permitem avaliar as duas abordagens. O preditivo apresenta custo médio 13,9\% maior e atendimento 4,11 pontos inferior à referência global; suas recomendações coincidem com e sem perfil explícito. Diagnósticos identificam poucos exemplos por classe, instabilidade dos rótulos e falhas de seleção distintas dos limites físicos do estoque. Uma regra exploratória de recuo ao perfil reduz o custo do preditivo em 12,2\%, mas não demonstra superioridade sobre as referências. Oráculos retrospectivos quantificam a distância ao melhor controle disponível e ao planejamento com informação perfeita. O estudo oferece uma arquitetura e uma avaliação rastreável, sem comprovar vantagem geral da contextualização. Amostragem retrospectiva, calibrações específicas por base e uso do piloto para desenvolvimento exigem confirmação independente.
\end{abstract}
\begin{keyword}
Inventory control \sep Algorithm selection \sep Demand heterogeneity \sep Meta-learning \sep Service level
\end{keyword}
\end{frontmatter}

\section{Introdução}
Uma rede varejista decide quando e quanto repor para muitas combinações de loja e produto. Vendas frequentes, longos intervalos sem demanda e picos ocasionais impõem necessidades diferentes. Essas diferenças motivam contextualizar a decisão, mas não demonstram que uma regra por grupo ou um modelo supervisionado supere uma política única bem calibrada. O benefício precisa aparecer no custo e no atendimento em períodos posteriores à escolha \cite{Axsater2006InventoryControl}.

Comparar controladores e selecionar entre eles são problemas distintos. Uma política pode ser competitiva em média e inadequada para algumas séries. Reciprocamente, a existência de vencedoras diferentes não garante que suas identidades sejam previsíveis. A seleção de algoritmos fornece a formulação pertinente: características de uma instância orientam a escolha de um método de um portfólio \cite{rice1976,Kotthoff2014AlgorithmSelection}. Aqui, a instância é a série loja--produto e a decisão é a família de política, com parâmetros calibrados no histórico da própria série.

Propomos o \textit{Adaptive Inventory Policy Engine} (AIPE) para organizar caracterização, calibração, seleção e avaliação. Seu componente de escolha, o \textit{Policy Selection Engine} (PSE), tem duas abordagens. O PSE por perfil recebe o perfil observado e consulta uma associação construída na validação. O PSE preditivo estima a política a partir das características da série, podendo receber o perfil como entrada adicional. A primeira compartilha uma recomendação no grupo; a segunda permite decisões distintas dentro dele. A aplicação pretendida reutiliza evidências de outras séries para orientar a família a calibrar. No experimento, todas as candidatas são ajustadas na série reservada para medir suas consequências e construir os oráculos; não foi mensurada economia computacional de implantação.

Investigamos três perguntas: o PSE por perfil acrescenta valor frente à seleção global? O preditivo melhora essa decisão e utiliza o perfil explícito? Quanto separa a recomendação dos melhores resultados retrospectivos? A avaliação combina dados proprietários da Bahia e um piloto público M5. As bases respondem a partes complementares dessas perguntas, sem reunir custos absolutos nem tratar o piloto como confirmação externa.

A contribuição é a arquitetura de seleção e sua avaliação operacional rastreável. Não propomos novas famílias de controladores, nem reivindicamos novidade na associação entre classes e políticas. Mostramos onde a contextualização altera o compromisso custo--serviço, onde o aprendizado falha e como separar erros de seleção de insuficiência do portfólio ou de limitações físicas. Os resultados negativos delimitam as condições em que a complexidade adicional ainda não se justifica.

\section{Trabalhos relacionados}
A classificação ADI--CV$^2$ organiza padrões de frequência e tamanho da demanda \cite{syntetos2005}. Entretanto, classes estatísticas não determinam sozinhas a melhor reposição: custos, prazo, estoque inicial e atendimento também importam. Mohammaditabar et al.\ \cite{Mohammaditabar2012Integrated} integram classificação e escolha de políticas; Sheikh-Zadeh et al.\ \cite{SheikhZadeh2021PBIC} propõem classificação orientada ao desempenho em sistemas de múltiplos itens e escalões. Esses antecedentes situam o AIPE como uma instanciação de seleção contextual, não como a primeira integração entre classificação e inventário. Seus métodos não foram reproduzidos como comparadores.

Na literatura da revista, Tadayonrad e Ndiaye \cite{sca_f81ad2b3e023} aproximam previsão e consequências operacionais. Theunissen et al.\ \cite{sca_c5c06dcbc9e5} discutem a informação decisória preservada pela classificação. Essa perspectiva orienta nossa avaliação: reconhecer um perfil ou acertar o nome de uma política só tem utilidade se melhorar a decisão. Não implementamos o método topológico desses autores.

Zabraoui et al.\ \cite{sca_27e994c54848} comparam algoritmos de inventário e combinações de busca genética e aprendizado por reforço com dados M5. A arquitetura GA--DQN integra nosso portfólio, com GA e DQN como controles, sem ser apresentada como contribuição algorítmica. Granularidade, filtros, custos e implementação local impedem uma reprodução exata de seu ranking.

Políticas limitadas de estoque-base oferecem outra referência para vendas perdidas \cite{Xin2021CappedBaseStock}. Suas propriedades teóricas não garantem desempenho para nossa calibração. Os resultados de reforço também dependem do estado, da recompensa e do orçamento \cite{Boute2022DRLInventoryRoadmap}. O AIPE separa essas implementações da escolha entre candidatas. A avaliação do seletor precisa, ainda, separar desenvolvimento e teste para evitar otimismo após escolher configurações \cite{CawleyTalbot2010}.

\section{AIPE e formulação da seleção}
\subsection{Da demanda à recomendação}
A Figura~\ref{fig:aipe} apresenta o fluxo. Cada histórico é preparado na unidade temporal da base e descrito por volume, dispersão, intermitência e estrutura temporal. As candidatas são calibradas no treino. Custo e serviço na validação sustentam a associação por perfil ou os rótulos do preditivo. A recomendação é avaliada no teste reservado.

\begin{figure}[htbp]\centering
\includegraphics[width=\linewidth]{figures/aipe_conceito_v2_ptbr.pdf}
\caption{AIPE e duas abordagens do PSE. O perfil pode orientar uma associação compartilhada ou entrar em um modelo que estima a política por série. Séries e associações A/B/C são ilustrativas. Ganhos são objeto de avaliação, não propriedade presumida da arquitetura.}\label{fig:aipe}
\end{figure}

A adaptação significa recalcular descritores e revisar a recomendação entre janelas. Não há troca de controlador a cada transação nem implantação online validada. Cada série constitui um sistema separado de vendas perdidas, sem estoque compartilhado, transbordo ou capacidade conjunta.

Seja $\Pi$ o portfólio, $x_i$ a descrição do treino e $g_i$ o perfil. O PSE por perfil recomenda $m(g_i)$; o preditivo recomenda $f(x_i,g_i)$ ou $f(x_i)$ na ablação. A seleção global usa uma mesma família para as séries da loja avaliada. Em todas as abordagens, os parâmetros são ajustados individualmente. A referência global, portanto, também utiliza o histórico de cada série.

\subsection{Custo, serviço e escolha}
Para demanda $d_t$, quantidade atendida $v_t$, saldo $I_t$ e pedido $q_t$, avaliamos
\begin{equation}
 C=\sum_t[hI_t+p(d_t-v_t)+K\mathbf1(q_t>0)+cq_t],
 \qquad F=\frac{\sum_t v_t}{\sum_t d_t}.
 \label{eq:metrics}
\end{equation}
O atendimento $F$ é \textit{fill rate}, não probabilidade de ciclo sem falta. Com demanda zero, permanece indefinido; o custo continua contabilizado. Usamos $h=1$, $p=5$, $K=50$ e $c=0{,}5$, em unidades do cenário, sem equivalência monetária calibrada.

Na validação $A$, a regra global minimiza $\bar C_{\pi,A}$ entre candidatas com $\bar F_{\pi,A}\geq\alpha$, com $\alpha=0{,}70$. O PSE por perfil aplica essa regra a cada grupo $A_g$. Casos sem demanda não entram na média de serviço. Sem candidata viável, escolhe-se maior serviço e menor custo, registrando contingência. Perfis sem suporte usam a global. Na Bahia, empate no custo viável prefere maior serviço e depois identificador; no M5, prefere identificador. Contingências empatadas em serviço preferem menor custo. Essas convenções são preservadas por base e registradas no código.

O rótulo supervisionado aplica a restrição à própria série de validação: menor custo viável ou, sem viável, maior serviço e menor custo. Se todas as opções têm serviço indefinido, utiliza menor custo e mantém a marcação de demanda zero. Classificador e seletores agregados não impõem a mesma distribuição da restrição. Incluímos a moda dos rótulos como controle para não atribuir toda diferença à capacidade preditiva.

\subsection{Características e perfis}
O preditivo recebe 17 descritores: ADI, CV$^2$, intermitência, proporção de zeros, maior e média das sequências de zeros, \textit{burstiness}, assimetria, curtose, entropia, tendência, autocorrelação, média, desvio-padrão, CV, extensão e número de ciclos positivos. O perfil categórico deriva desses descritores, sem constituir informação independente.

As regras são ordenadas (Tabela~\ref{tab:profiles}). ADI é $n/\max(n_+,1)$ e CV$^2$ usa a dispersão de todos os ciclos, incluindo zeros, diferindo da taxonomia original baseada em tamanhos positivos. Os nomes são descritores heurísticos, sem comprovar sazonalidade ou impacto econômico.

\begin{table}[htbp]\centering\small
\caption{Primeira condição satisfeita determina o perfil; casos residuais recebem o perfil esparso. Cortes fixos, sem ajuste no teste.}\label{tab:profiles}
\begin{tabular}{@{}p{3.5cm}p{7.4cm}@{}}\toprule
Perfil & Condição\\\midrule
Esparso & ADI $\geq1{,}32$, CV$^2\geq0{,}49$ e sequência de zeros $\geq5$\\
Volatilidade / sazonalidade & Autocorrelação na defasagem 17 $\geq0{,}25$ ou burstiness $\geq0{,}30$\\
Tendência instável & Tendência normalizada absoluta $\geq0{,}05$ ou entropia $\geq2$\\
Baixo volume estável & Média $\leq2$ e CV$^2<0{,}49$\\
Venda frequente & ADI $<1{,}32$\\\bottomrule
\end{tabular}\end{table}

Burstiness compara o desvio e a média dos intervalos positivos pela razão de sua diferença e soma, com regularização numérica. Tendência é a inclinação OLS dividida pela média; entropia usa dez classes de mesma largura. A autocorrelação na defasagem 17 recebe zero com até 18 ciclos ou variância desprezível. Essa convenção denota estimativa indisponível, não ausência de sazonalidade; com 23 ciclos, há somente seis pares.

\section{Desenho experimental}
\subsection{Bases e unidades de análise}
A Tabela~\ref{tab:datasets} distingue as avaliações. Na Bahia, transações complementares são somadas por loja, produto e ciclo; ciclos sem venda recebem zero. A coorte retrospectiva contém 145 séries, 115 lojas e nove produtos. Os 38 ciclos disponíveis, de 202301 a 202504, são unidades comerciais da fonte, não meses.

O M5 é uma fonte pública de vendas associada a uma competição de previsão \cite{Makridakis2022M5Accuracy}. O piloto sorteia 30 séries FOODS, três por loja nas dez lojas. Na primeira origem, o volume médio do treino define três estratos por loja, com uma escolha em cada estrato e semente 20261015. As séries são mantidas nas três origens. Usamos ciclos completos de 21 dias, sem aparar picos por IQR.

\begin{table}[htbp]\centering\small
\caption{Matrizes completas analisadas. Casos série--origem e sementes não são séries independentes.}\label{tab:datasets}
\begin{tabular}{@{}lrr@{}}\toprule
Dimensão & Bahia & Piloto M5\\\midrule
Séries / lojas & 145 / 115 & 30 / 10\\
Origens de treino (ciclos) & 17 e 23 & 50, 64 e 78\\
Validação / teste por origem & 7 / 7 & 7 / 7\\
Casos série--origem & 290 & 90\\
Ajustes de políticas & 4.350 & 1.350\\
Testes sem demanda & 7 & 4\\
Seletores & Global e por perfil & Global, por perfil e preditivo\\\bottomrule
\end{tabular}\end{table}

A elegibilidade exige três ciclos positivos no prefixo. O sorteio M5 usa 885 séries elegíveis na primeira origem de uma coorte histórica de 2.408 identificadores. Não representa todas as FOODS nem séries elegíveis somente depois. O recorte é inspirado no contexto de Zabraoui et al., sem afirmar identidade com sua amostra. A rodada maior interrompida não fornece conclusões deste artigo.

\subsection{Simulador, portfólio e calibração}
O prazo exato é de dois ciclos. A política observa o estado anterior ao recebimento; entregas chegam antes da demanda. Quantidades podem ser fracionárias. Faltas são vendas perdidas, sem pedidos pendentes. Manutenção incide no saldo após demanda e o pedido é custeado na emissão, sem crédito terminal.

Validação e teste começam em estados comuns às candidatas, produzidos por uma referência EOQ: seis ciclos passados aquecem o sistema, e sua trajetória na validação determina estoque e pipeline do teste. Esses estados usam o passado, mas não representam a continuação particular de cada candidata. Escalas são calculadas no treino.

O portfólio contém EOQCycle, sSCycle, CappedBaseStock, GA, PPO, GA--DQN e DQN. EOQCycle usa lote $\max(1,\sqrt{2\mu K/h})$ e ponto de pedido $2\mu+1{,}28\sigma\sqrt2$; sSCycle repõe até o ponto acrescido do lote. CappedBaseStock limita a quantidade da reposição até um alvo e calibra alvo e teto em grade $20\times12$. GA ajusta gatilho, lote e proteção com população 100 e 50 gerações. DQN usa Double DQN dueling; PPO usa otimização proximal. Ambos têm duas camadas de 64 unidades, 20 ações e no mínimo 20 mil transições, arredondadas ao episódio completo. GA--DQN acrescenta busca genética e mil demonstrações, mantendo gatilho e piso iguais nas demonstrações, treino e inferência.

Os quatro métodos estocásticos usam sementes 11, 29 e 47; as três referências determinísticas têm um ajuste. As sementes são promediadas antes da avaliação por caso. O híbrido tem esforço superior ao DQN, sem equivalência de orçamento. A calibração pode penalizar déficit de serviço. Na Bahia, CappedBaseStock também penaliza frequência de ruptura e amplificação dos pedidos com pesos unitários; no M5, ambos os pesos são zero. Essa diferença impede atribuir contrastes entre bases somente à demanda. Cada matriz e código arquivado identificam a instanciação avaliada.

\subsection{Avaliação e desenvolvimento do piloto}
Em cada origem, uma loja inteira fica fora do conjunto que define a escolha. O PSE usa a validação das demais lojas e recomenda para o teste da loja reservada. Controladores continuam calibrados no histórico da própria série: avaliamos transferência do seletor com histórico disponível, não implantação sem dados prévios.

No M5, XGBoost tem 200 árvores, profundidade três, taxa 0,05, sem subamostragem e com semente fixa. Imputação, codificação categórica e classes são ajustadas dentro de cada fold. Classes raras não são removidas; treinos de classe única usam classificador constante. Cada treinamento reúne 27 exemplos. A ablação sem perfil mantém os demais componentes. Não se usam origens posteriores para recomendar em uma origem anterior.

Após o diagnóstico, examinamos uma única extensão: recuar ao PSE por perfil se o maior escore de classe for inferior a 0,60 ou a classe prevista tiver menos de três exemplos de treino. Não há busca de limiares no teste, mas a motivação veio dos resultados observados. Seu desempenho é exploratório, e o piloto passa a ser conjunto de desenvolvimento \cite{CawleyTalbot2010}. O escore não é probabilidade calibrada de viabilidade.

Na Bahia, os testes compartilham um ciclo; no M5, ocupam os ciclos 58--64, 72--78 e 86--92, sem sobreposição. Origens repetem séries. Intervalos exploratórios usam bootstrap pareado por loja, preservando séries e origens: 5.000 sorteios na Bahia e 2.000 no M5. Seletores e controladores permanecem fixos; a incerteza é condicional, sem novos treinamentos. Dez lojas do piloto limitam sua precisão.

\subsection{Oráculos e informação perfeita}
O melhor controlador único no teste é um diagnóstico retrospectivo. O oráculo do portfólio escolhe, por série--origem, a candidata de menor custo viável nas médias de teste entre sementes. Essas escolhas não fornecem rótulos ao treinamento. O acerto exige serviço mínimo e empate numérico no menor custo; concordância nominal é informada separadamente.

No M5, um programa linear inteiro misto conhece toda a demanda de teste e otimiza pedidos sob os mesmos custos, prazo e condições iniciais. Quantidades são contínuas; variáveis binárias representam emissão e complementaridade entre saldo e falta. É um ótimo determinístico condicional, não um controlador estocástico implantável. Se 0,70 for fisicamente impossível antes das primeiras entregas, maximiza atendimento e depois minimiza custo. Soluções são reexecutadas no simulador e somente as comprovadamente ótimas recebem essa denominação.

A eficiência de custo é $100C^*/C_{\mathrm{PSE}}$, apenas quando referência e recomendação satisfazem o mesmo critério de serviço; ambos os custos zero equivalem a 100\%. Serviço insuficiente não recebe crédito por custo baixo. A comparação com o MILP usa cada realização de treinamento viável: uma média entre sementes pode satisfazer o alvo sem que todas o satisfaçam. Repetir políticas determinísticas nessa análise não gera novas observações independentes.

\input{aipe_bahia_resultados}
\input{m5_pse_20261015}

\section{Discussão}
O AIPE explicita quando vale mudar a família de política entre séries. Na Bahia, o PSE por perfil troca custo por atendimento e diferencia escolhas em somente uma origem. No piloto M5, a global oferece compromisso melhor que o preditivo, enquanto o PSE por perfil permanece próximo da referência. Não há evidência de superioridade geral da contextualização.

A dificuldade do preditivo não se reduz ao algoritmo. Há sete classes e 27 exemplos por fold; rótulos variam com período e semente. Um classificador pode reproduzir o treino e falhar na loja reservada. O alvo categórico também descarta quão próximas as candidatas estão em custo e serviço. Isso motiva investigar alvos mais estáveis e critérios sensíveis à perda operacional, sem escolher soluções pelos mesmos testes usados para justificá-las.

O recuo limita recomendações individuais pouco sustentadas e aproxima seu custo ao da global, sem atendimento superior. Não resolve falta de opções viáveis nem faltas inevitáveis. A decomposição indica se a próxima intervenção deve mudar seletor, calibração ou condições de abastecimento.

A concentração de perfis limita sua avaliação. Na primeira origem M5, 826 das 885 séries elegíveis são esparsas, 57 têm tendência e apenas duas pertencem ao grupo volatilidade/sazonalidade. A baixa diversidade não decorre somente do sorteio. Aumentar a amostra sem registrar suporte por grupo não assegura avaliação representativa de todos os perfis.

Meta média, atendimento ponderado por unidades e garantia por série são exigências diferentes. Média de 0,70 pode coexistir com numerosas violações individuais. Essa definição precisa anteceder a seleção; modificá-la para favorecer o resultado mudaria o problema. Os custos e o limiar deste cenário não constituem recomendação comercial.

\section{Limitações e próximos testes}
Vendas podem omitir demanda censurada por falta. Coortes históricas, elegibilidade na primeira origem e horizontes de sete ciclos restringem população e período. Produtos compartilhados entre lojas podem gerar dependência não capturada inteiramente pelo bootstrap.

Não se trata de experimento causal de implantação. Custos não são financeiramente calibrados, não há suprimento conjunto e o estado inicial comum não reproduz trajetórias particulares. A calibração distinta da CappedBaseStock e os esforços desiguais limitam comparações entre algoritmos e bases.

A extensão com recuo foi desenvolvida no piloto; suas métricas não confirmam independentemente a melhoria. A próxima avaliação exige mais exemplos por classe, suporte para perfis raros, restrições de serviço comparáveis e validação interna por loja para ajustes do seletor. Séries reservadas devem ficar fora dessas decisões. Harmonização da calibração e sensibilidades de custo, prazo e horizonte precedem conclusões mais amplas.

\section{Conclusão}
Apresentamos o AIPE com PSE por perfil e preditivo, avaliando recomendações separadamente dos controladores. Na Bahia, o perfil acrescenta 3,89 pontos de atendimento com aumento pontual de custo de 1,93\%. No piloto M5, o preditivo custa 13,9\% mais e atende 4,11 pontos menos que a global. O recuo melhora esse preditivo, sem estabelecer vantagem sobre as referências.

Heterogeneidade e portfólio diverso não bastam para tornar a escolha aprendível. Oráculos e decomposição das falhas distinguem distância ao melhor candidato de restrições que o seletor não supera. A proposta permanece em desenvolvimento; avaliação independente deverá demonstrar em quais contextos sua complexidade oferece benefício operacional.

\section*{Disponibilidade de dados e código}
Os dados brasileiros são proprietários e não são redistribuídos; o M5 é público. Código, configurações e fontes estão em \url{https://github.com/matheus013/msc}. Dados, bancos e modelos ficam fora da versão atual do repositório. Manifestos e snapshots locais identificam as matrizes; pacote arquivado e derivados autorizados deverão acompanhar a submissão. Não se promete acesso aos registros restritos mediante solicitação.

\section*{Autoria, financiamento e interesses}
Autoria, contribuições CRediT, financiamento, agradecimentos e interesses devem ser confirmados pelos autores responsáveis. O manuscrito em português requer tradução, revisão humana e fechamento das declarações antes da submissão.

\section*{Uso de inteligência artificial na preparação do manuscrito}
OpenAI Codex auxiliou organização bibliográfica, leitura de código, análise, visualização e redação. A figura conceitual usa desenhos vetoriais programáticos; os gráficos quantitativos derivam das matrizes identificadas. Conferência e responsabilidade pelo texto final cabem aos autores humanos. Sistemas de IA não são autores.

\bibliographystyle{elsarticle-num}
\bibliography{references_ptbr}
\end{document}
